\section{Экспериментальный протокол}

В этом разделе описывается протокол, по которому были получены сравнимые результаты для MTP, MACE и MatterSim. Здесь не анализируется, какая модель оказалась лучше на конкретной системе; цель состоит в том, чтобы зафиксировать данные, преобразования форматов, режимы обучения и правила расчета итоговых метрик.

Общая схема эксперимента была одинаковой для всех систем. Атомистические конфигурации приводились к форматам используемых инструментов, после чего для каждой целевой системы обучались MTP-потенциалы с нуля и рассматривались универсальные предобученные модели MatterSim и MACE. Для универсальных моделей основной режим состоял в тонкой настройке на целевых данных; оценки без дообучения использовались только как диагностическая часть протокола при наличии проверяемой цепочки источников.

\subsection{Данные и контрольные разбиения}

Основная линия экспериментов включает три атомистические системы: organic-датасет, воду H2O и многокомпонентный сплав MoNbTaVW. Они отличаются как химическим составом, так и смыслом контрольного разбиения. Поэтому в работе не используется единая интерпретация всех разбиений: для organic-датасета основной интерес представляет перенос качества между температурными режимами, а для H2O и MoNbTaVW анализ проводится на отложенном валидационном разбиении.

\begin{table}[h]
\centering
\small
\caption{Системы и контрольные разбиения в основной линии экспериментов}
\label{tab:protocol_datasets_splits}
\begin{tabularx}{\textwidth}{@{}l>{\raggedright\arraybackslash}X>{\raggedright\arraybackslash}X>{\raggedright\arraybackslash}X@{}}
\toprule
Система & Обучающие данные & Контрольные разбиения & Смысл оценки \\
\midrule
organic-датасет & данные, связанные с 300K & test\_300K, test\_600K, test\_1200K & качество при переносе между температурными режимами \\
H2O & train-разбиение & valid & качество на отложенных конфигурациях \\
MoNbTaVW & train-разбиение & valid & качество на отложенных конфигурациях \\
\bottomrule
\end{tabularx}
\end{table}

Для organic-датасета разбиение test\_300K близко к обучающему температурному режиму, а test\_600K и test\_1200K используются как более сложные контрольные режимы. Поэтому эти результаты интерпретируются как проверка качества на новых конфигурациях и устойчивости при изменении температуры. Для H2O и MoNbTaVW температурный перенос в основной линии результатов не рассматривается: там качество оценивается на валидационном разбиении.

Метрики на обучающих разбиениях использовались только как диагностический показатель и не являются основой для итоговых выводов. Основные выводы формулируются по контрольным разбиениям: test\_300K, test\_600K и test\_1200K для organic-датасета и валидационным разбиениям для H2O и MoNbTaVW.

\subsection{Подготовка данных и преобразование форматов}

Для сравнения MTP, MACE и MatterSim одни и те же конфигурации требовалось использовать в разных программных инструментах. В проекте применялись представления на основе \texttt{ASE/extxyz}, формат \texttt{.cfg}, используемый в MLIP, и JSON-представление, совместимое с обработкой в MLIP-4. Поэтому отдельной частью работы стала подготовка преобразований между форматами.

При конвертации сохранялись величины, необходимые для обучения и оценки: химические типы атомов, координаты, параметры ячейки для периодических систем, энергии конфигураций и силы. Нарушение этой цепочки могло бы изменить MAE и RMSE либо повлиять на построение соседских окружений в периодических системах.

Скрипты преобразования сохранены в части репозитория, отвечающей за происхождение материалов. Они не являются основной точкой запуска очищенного репозитория, но фиксируют, как обеспечивалась совместимость данных между MTP, MACE и MatterSim.

\subsection{Режимы обучения и оценки моделей}

В работе сравниваются две стратегии построения межатомного потенциала. Первая стратегия состоит в обучении специализированной модели с нуля на данных целевой системы. В этой роли используются MTP-потенциалы уровней MTP-16 и MTP-20. Такая модель не использует внешнее предобучение и получает информацию только из данных конкретной системы.

Вторая стратегия состоит в использовании универсального предобученного потенциала и его последующей адаптации к целевой системе. В этой роли рассматриваются MatterSim и MACE. Для них предобученные параметры служат начальной точкой оптимизации, а тонкая настройка проводится на целевых данных. Такое сравнение позволяет сопоставить обучение с нуля и перенос информации из внешнего предобучения.

\begin{table}[h]
\centering
\small
\caption{Основные режимы, сравниваемые в экспериментальном протоколе}
\label{tab:protocol_training_modes}
\begin{tabularx}{\textwidth}{@{}l>{\raggedright\arraybackslash}X>{\raggedright\arraybackslash}X@{}}
\toprule
Модель или протокол & Режим & Роль в сравнении \\
\midrule
MTP-16, MTP-20 & обучение с нуля & специализированная базовая линия \\
MatterSim & тонкая настройка предобученной модели & универсальный потенциал после адаптации \\
MACE default & тонкая настройка предобученной модели & универсальный потенциал после адаптации \\
MACE-EW50 & тонкая настройка с увеличенным весом ошибки энергии & отдельный протокол MACE для MoNbTaVW \\
\bottomrule
\end{tabularx}
\end{table}

Для MTP сравнивались два уровня сложности: MTP-16 и MTP-20. Уровень задает размер и выразительность базиса, поэтому MTP-20 рассматривается как более сложная специализированная базовая линия. Сведения о числе базисных функций, параметрах и дополнительных проверках вынесены в приложение.

Для MatterSim и MACE основной интерес представляет качество после тонкой настройки. Оценки предобученной модели без адаптации используются только как вспомогательный элемент протокола и не подменяют основные строки результатов.

Для MoNbTaVW отдельно фиксируется различие между MACE default и MACE-EW50. Эти строки нельзя рассматривать как дубли одного и того же запуска: MACE-EW50 соответствует отдельному протоколу тонкой настройки, в котором увеличивается относительный вес ошибки энергии. Поэтому в результатах он интерпретируется как отдельный экспериментальный режим, а не как альтернативное название модели MACE.

\subsection{Расчет метрик и границы воспроизводимости}

После обучения или тонкой настройки модели оценивались на соответствующих контрольных разбиениях. Для каждой модели рассчитывались MAE и RMSE по энергии в meV/atom, а также MAE и RMSE по силам в meV/Å. Формулы метрик приведены в разделе 2; здесь фиксируется только единая система единиц и структура итоговых таблиц.

Итоговые численные результаты не выбирались заново при написании отчета. Для основной линии использовались проверенные таблицы и сохраненные результаты оценки, чтобы не смешивать старые промежуточные запуски, презентационные сопроводительные таблицы и финальные строки отчета. При интерпретации учитываются только строки с проверяемой связью между моделью, системой, разбиением, метриками и исходной оценкой.

Границы воспроизводимости также фиксируются явно. Очищенный репозиторий позволяет проверить структуру результатов, скрипты агрегации, небольшие сохраненные оценки и цепочку происхождения. Он не предназначен для полного переобучения всех моделей с нуля: исходные датасеты, веса моделей, контрольные точки моделей и тяжелые промежуточные артефакты не включены физически. Поэтому воспроизводимость в рамках отчета понимается как проверяемость итоговых таблиц, протокола и цепочки происхождения, а не как полный повтор всех вычислительных экспериментов из одного очищенного репозитория.
